% ===================================================================
% appendices/A_runs.tex -- Catalogue of all simulation runs (V1: 52 launches), with V0 and Astra history.
% ASCII only. The table itself (A_runs_table.tex) is generated by tools/make_appendix_assets.py
% from results/RUN_INDEX.csv and the execution.json of every run.
% ===================================================================
\chapter{Catalogue of all simulation runs}\label{ch:app-runs}

This appendix lists every launch of the V1 package, including those that failed or were retracted, so that each
simulated number of the report can be traced to its run folder and so that no run can silently disappear from the
record. The catalogue (\cref{tab:run-catalogue}) is generated from \file{results/RUN_INDEX.csv} and the
\file{execution.json} of each run by \file{tools/make_appendix_assets.py}; no value in it is typed by hand. The run
identifier \run{NNNN} names the folder \file{results/runs/run_NNNN_<film>_<label>}. The overrides column states how the
run differs from the reference configuration of \file{config/iwo_material_model.yaml}, \file{config/geometry.yaml} and
\file{config/solver.yaml} (dotted paths as accepted by \file{scripts/build_iwo_decks.py}); metric definitions are those
of \cref{sec:der-metrics}, and the numerical settings of \cref{ch:numerics}.

\section{Composition of the 52 launches}\label{sec:app-runs-composition}

\begin{table}[htbp]
\centering
\small
\caption{Runs by purpose (categories of \protect\file{results/RUN_INDEX.csv}).}
\label{tab:app-runs-categories}
\begin{tabularx}{\linewidth}{@{}lrX@{}}
\toprule
Category & Runs & Identifiers \\
\midrule
Baseline thickness laws & 2 & \run{0001} (2 nm), \run{0002} (13.2 nm) \\
Calibration stages & 7 & \run{0003}, \run{0004}, \run{0006}, \run{0007}, \run{0008}, \run{0009}, \run{0011} \\
Held-out validation of \SI{6.3}{nm} & 2 & \run{0005}, \run{0014} \\
Final calibrated, DOS 96/48 & 2 & \run{0012} (2 nm), \run{0013} (13.2 nm) \\
Device-specific tuning of \SI{6.3}{nm} & 1 & \run{0015} \\
Counterfactual without confinement & 2 & \run{0016} (2 nm), \run{0017} (13.2 nm) \\
One-at-a-time sensitivity (2 nm) & 8 & \run{0020}--\run{0027} \\
Numerical checks & 5 & \run{0018}, \run{0019}, \run{0028}, \run{0029}, \run{0036} \\
\SI{6.3}{nm} hypotheses & 9 & \run{0010}, \run{0030}--\run{0035}, \run{0037}, \run{0052} \\
DOS 384/192 recalibration & 3 & \run{0038}, \run{0039}, \run{0040} \\
Predictions, $\Id$--$\Vd$ & 3 & \run{0041}--\run{0043} \\
Predictions, elevated temperature & 4 & \run{0044}--\run{0047} \\
Isolation tests (gate WF, $\mstar$) & 4 & \run{0048}--\run{0051} \\
\midrule
Total & 52 & \\
\bottomrule
\end{tabularx}
\end{table}

The record holds 52 recorded launches, of which 49 produced scored currents: 46 carry the status
\texttt{REAL\_ATLAS\_SCORED} and three \texttt{REAL\_ATLAS\_PREDICTION} ($\Id$--$\Vd$ families without a measurement to
score against), while three produced no usable curve (\run{0010}, \run{0011}, \run{0018}); one scored run (\run{0027}) is retracted. The runs \run{0001}--\run{0009}
use the reduced structure and the pointwise sweep; every run from \run{0012} on uses the physical structure and the
stepped sweep (\cref{sec:num-sweep}). DOS 384/192 is used by the check \run{0029} and by all runs from \run{0038} to
\run{0052}; all other runs, including the pre-registered B0 test \run{0037}, use 96/48.

% generated by tools/make_appendix_assets.py from digest/runs.csv (execution.json + comparison.csv of each run); do not edit by hand
{\scriptsize\setlength{\tabcolsep}{2.5pt}
\begin{longtable}{@{}l r p{3.3cm} p{2.3cm} r r r r r r r@{}}
\caption{Catalogue of all V1 ATLAS launches and the measured references. $\SSfix{-11}{-10}$: fixed-current slope between $10^{-11}$ and $10^{-10}$\Aum{} (mV/dec); RMSE: active-region log10 error (dec). Metrics of elevated-temperature runs are compared with the 300\,K data only for reference.}\label{tab:run-catalogue}\\
\toprule run & $t$ (nm) & label & category & $\Vthcc$ (V) & $\Vthlin$ (V) & $\SSfix{-11}{-10}$ & $\SScc$ & $\muFE$ & $\Ion$ (A/$\mu$m) & RMSE \\ \midrule \endfirsthead
\toprule run & $t$ (nm) & label & category & $\Vthcc$ & $\Vthlin$ & $\SSfix{-11}{-10}$ & $\SScc$ & $\muFE$ & $\Ion$ & RMSE \\ \midrule \endhead
\bottomrule \endfoot
\run{0001} & 2 & v1\_baseline\_laws & baseline laws & 0.995 & 1.838 & 125.3 & -- & 10.19 & 3.71e-07 & 3.928 \\
\run{0002} & 13.2 & v1\_baseline\_laws & baseline laws & 0.369 & 1.285 & 92.4 & -- & 44.97 & 2.42e-06 & 1.030 \\
\run{0003} & 2 & v1\_stage3\_qf1p73e12 & calibration stage & 0.686 & 1.561 & 125.0 & -- & 10.46 & 4.72e-07 & 0.056 \\
\run{0004} & 13.2 & v1\_stage3\_qf1p73e12 & calibration stage & 0.059 & 1.001 & 92.0 & 195.4 & 45.60 & 2.86e-06 & 0.086 \\
\run{0005} & 6.3 & v1\_stage5\_shared\_qf\_prediction & validation (held-out 6.3 nm) & 0.350 & 1.215 & 122.3 & 284.9 & 9.14 & 5.11e-07 & 0.678 \\
\run{0006} & 31.8 & v1\_stage5\_shared\_qf & calibration stage & -2.710 & -0.168 & 84.9 & 210.2 & 39.90 & 3.38e-06 & 1.050 \\
\run{0007} & 6.3 & v1\_stage5b\_device\_qf8p7e10 & calibration stage & 0.644 & 1.474 & 122.8 & 284.7 & 8.95 & 4.28e-07 & 0.082 \\
\run{0008} & 2 & v1\_stage6\_mu18p13 & calibration stage & 0.676 & 1.561 & 122.5 & 289.1 & 11.29 & 5.09e-07 & 0.037 \\
\run{0009} & 13.2 & v1\_stage6\_mu61p9 & calibration stage & 0.053 & 1.001 & 90.8 & 192.5 & 49.00 & 3.07e-06 & 0.081 \\
\run{0010} & 31.8 & v1\_stage6\_backsheet\_hyp\_A & 6.3 nm hypothesis & -- & -- & -- & -- & -- & -- & -- \\
\multicolumn{11}{@{}l}{\hspace{1em}\textit{aborted in SOLVE INIT, no currents}} \\
\run{0011} & 6.3 & v1\_stage6\_mu11p69 & calibration stage & -- & -- & -- & -- & -- & -- & -- \\
\multicolumn{11}{@{}l}{\hspace{1em}\textit{stopped with the chain; no currents}} \\
\run{0012} & 2 & v2\_full\_stepped\_final & final calibrated (DOS 96/48) & 0.676 & 1.559 & 122.5 & 289.1 & 11.27 & 5.09e-07 & 0.037 \\
\run{0013} & 13.2 & v2\_full\_stepped\_final & final calibrated (DOS 96/48) & 0.053 & 0.999 & 90.8 & 192.5 & 48.95 & 3.07e-06 & 0.081 \\
\run{0014} & 6.3 & v2\_full\_stepped\_validation\_shared\_params & validation (held-out 6.3 nm) & 0.350 & 1.212 & 122.3 & 284.9 & 9.13 & 5.11e-07 & 0.678 \\
\run{0015} & 6.3 & v2\_full\_stepped\_tuned\_qf8p7e10\_mu11p69 & device-specific tuning & 0.651 & 1.471 & 124.1 & 288.8 & 8.42 & 4.03e-07 & 0.063 \\
\run{0016} & 2 & sens\_no\_confinement & counterfactual (no confinement) & 0.361 & 1.260 & 132.1 & 303.7 & 11.33 & 6.18e-07 & 0.846 \\
\run{0017} & 13.2 & sens\_no\_confinement & counterfactual (no confinement) & 0.030 & 0.979 & 92.0 & 193.3 & 48.92 & 3.1e-06 & 0.114 \\
\run{0018} & 13.2 & check\_A\_mesh0p7 & numerical check & -- & -- & -- & -- & -- & -- & -- \\
\multicolumn{11}{@{}l}{\hspace{1em}\textit{timed out at 900 s; rerun as run\_0028}} \\
\run{0019} & 2 & check\_C\_dos192\_96 & numerical check & 0.677 & 1.541 & 122.7 & 289.3 & 11.35 & 5.19e-07 & 0.040 \\
\run{0020} & 2 & sens\_chi\_m0p1 & one-at-a-time sensitivity & 0.776 & 1.649 & 122.5 & 289.1 & 11.19 & 4.74e-07 & 0.463 \\
\run{0021} & 2 & sens\_chi\_p0p1 & one-at-a-time sensitivity & 0.576 & 1.467 & 122.5 & 289.1 & 11.34 & 5.45e-07 & 0.361 \\
\run{0022} & 2 & sens\_nd\_x2 & one-at-a-time sensitivity & 0.667 & 1.550 & 121.5 & 289.8 & 11.28 & 5.13e-07 & 0.054 \\
\run{0023} & 2 & sens\_nt\_x1p5 & one-at-a-time sensitivity & 0.758 & 1.841 & 144.4 & 357.8 & 10.50 & 3.82e-07 & 0.263 \\
\run{0024} & 2 & sens\_wta\_p5meV & one-at-a-time sensitivity & 0.696 & 1.557 & 130.0 & 290.0 & 11.30 & 5.11e-07 & 0.063 \\
\run{0025} & 2 & sens\_dit\_x3 & one-at-a-time sensitivity & 0.701 & 1.582 & 123.9 & 289.7 & 11.25 & 5e-07 & 0.098 \\
\run{0026} & 2 & sens\_eps\_10p55 & one-at-a-time sensitivity & 0.675 & 1.558 & 122.1 & 289.3 & 11.34 & 5.13e-07 & 0.037 \\
\run{0027} & 2 & sens\_overlap\_4um & one-at-a-time sensitivity & 0.675 & 1.564 & 122.2 & 290.0 & 11.07 & 4.98e-07 & 0.038 \\
\multicolumn{11}{@{}l}{\hspace{1em}\textit{MALFORMED: x.mesh ends at 24 um, drain collapsed; overlap result retracted (S5)}} \\
\run{0028} & 13.2 & check\_A\_mesh0p7\_t1500 & numerical check & 0.053 & 0.999 & 90.7 & 192.5 & 48.94 & 3.07e-06 & 0.082 \\
\run{0029} & 2 & check\_C\_dos384\_192 & numerical check & 0.677 & 1.537 & 122.7 & 289.3 & 11.37 & 5.22e-07 & 0.041 \\
\run{0030} & 6.3 & hyp\_H3\_backQ\_m1p64e12 & 6.3 nm hypothesis & 0.627 & 1.352 & 102.7 & 231.3 & 8.93 & 4.61e-07 & 0.237 \\
\run{0031} & 6.3 & hyp\_H1\_tphys\_5p3 & 6.3 nm hypothesis & 0.386 & 1.265 & 124.5 & 290.5 & 8.96 & 4.87e-07 & 0.613 \\
\run{0032} & 6.3 & hyp\_H2\_Nd0\_1e16 & 6.3 nm hypothesis & 0.380 & 1.227 & 122.5 & 280.7 & 9.03 & 5.02e-07 & 0.634 \\
\run{0033} & 6.3 & hyp\_H4\_front\_Dit\_x5 & 6.3 nm hypothesis & 0.401 & 1.258 & 127.3 & 286.3 & 9.09 & 4.97e-07 & 0.595 \\
\run{0034} & 6.3 & hyp\_H6\_noQC\_Qf\_from\_2nm & 6.3 nm hypothesis & 0.587 & 1.426 & 125.1 & 287.7 & 8.93 & 4.4e-07 & 0.211 \\
\run{0035} & 6.3 & hyp\_H5\_combo\_t5p8\_Nd1e16\_backQ\_m1e12 & 6.3 nm hypothesis & 0.565 & 1.336 & 109.6 & 249.5 & 8.81 & 4.59e-07 & 0.246 \\
\run{0036} & 6.3 & check\_L\_mesh0p7\_tuned & numerical check & 0.651 & 1.471 & 124.2 & 288.9 & 8.42 & 4.03e-07 & 0.062 \\
\run{0037} & 6.3 & hyp\_S2\_repartition\_Nt2p66e19\_mu19p6 & 6.3 nm hypothesis & 0.654 & 1.856 & 173.3 & 410.9 & 10.40 & 3.73e-07 & 0.270 \\
\run{0038} & 2 & recal\_dos384\_2p0 & DOS 384/192 recalibration & 0.680 & 1.537 & 123.5 & 290.9 & 11.10 & 5.09e-07 & 0.043 \\
\run{0039} & 6.3 & recal\_dos384\_6p3\_tuned & DOS 384/192 recalibration & 0.655 & 1.466 & 124.8 & 290.7 & 8.27 & 3.97e-07 & 0.059 \\
\run{0040} & 13.2 & recal\_dos384\_13p2 & DOS 384/192 recalibration & 0.055 & 0.995 & 91.2 & 193.8 & 47.89 & 3.01e-06 & 0.081 \\
\run{0041} & 2 & pred\_idvd\_2p0 & prediction ID-VD & -- & -- & -- & -- & -- & -- & -- \\
\multicolumn{11}{@{}l}{\hspace{1em}\textit{ID-VD prediction; output\_characteristics.csv}} \\
\run{0042} & 6.3 & pred\_idvd\_6p3\_tuned & prediction ID-VD & -- & -- & -- & -- & -- & -- & -- \\
\multicolumn{11}{@{}l}{\hspace{1em}\textit{ID-VD prediction; output\_characteristics.csv}} \\
\run{0043} & 13.2 & pred\_idvd\_13p2 & prediction ID-VD & -- & -- & -- & -- & -- & -- & -- \\
\multicolumn{11}{@{}l}{\hspace{1em}\textit{ID-VD prediction; output\_characteristics.csv}} \\
\run{0044} & 2 & pred\_T338K\_2p0 & prediction temperature & 0.653 & 1.519 & 123.5 & 300.2 & 9.22 & 4.28e-07 & 0.131 \\
\multicolumn{11}{@{}l}{\hspace{1em}\textit{elevated T; ATLAS tmu=1.5 applied (P-phonon as simulated)}} \\
\run{0045} & 2 & pred\_T358K\_2p0 & prediction temperature & 0.640 & 1.510 & 125.0 & 304.5 & 8.44 & 3.94e-07 & 0.186 \\
\multicolumn{11}{@{}l}{\hspace{1em}\textit{elevated T; ATLAS tmu=1.5 applied (P-phonon as simulated)}} \\
\run{0046} & 6.3 & pred\_T358K\_6p3 & prediction temperature & 0.606 & 1.432 & 123.7 & 307.3 & 6.43 & 3.16e-07 & 0.133 \\
\multicolumn{11}{@{}l}{\hspace{1em}\textit{elevated T; ATLAS tmu=1.5 applied (P-phonon as simulated)}} \\
\run{0047} & 13.2 & pred\_T358K\_13p2 & prediction temperature & -0.003 & 0.964 & 93.0 & 195.2 & 37.46 & 2.39e-06 & 0.182 \\
\multicolumn{11}{@{}l}{\hspace{1em}\textit{elevated T; ATLAS tmu=1.5 applied (P-phonon as simulated)}} \\
\run{0048} & 2 & sens\_wf\_p0p1\_2p0\_dos384 & isolation (WF) & 0.780 & 1.628 & 123.5 & 290.9 & 11.03 & 4.74e-07 & 0.470 \\
\run{0049} & 13.2 & sens\_wf\_p0p1\_13p2\_dos384 & isolation (WF) & 0.155 & 1.087 & 91.2 & 193.8 & 47.69 & 2.86e-06 & 0.221 \\
\run{0050} & 2 & sens\_mstar\_x1p3\_2p0\_dos384 & isolation (m*) & 0.637 & 1.494 & 113.3 & 273.2 & 11.23 & 5.3e-07 & 0.087 \\
\run{0051} & 13.2 & sens\_mstar\_x1p3\_13p2\_dos384 & isolation (m*) & 0.027 & 0.934 & 86.5 & 180.5 & 49.17 & 3.18e-06 & 0.122 \\
\run{0052} & 6.3 & hyp\_S4\_nearEc\_band\_2e12\_mu15p4 & 6.3 nm hypothesis & 0.648 & 1.587 & 125.9 & 304.4 & 10.56 & 4.68e-07 & 0.081 \\
\run{0053} & 2 & check\_overlap\_4um\_meshfixed\_dos384 & numerical check & -- & -- & -- & -- & -- & -- & -- \\
\run{0054} & 2 & check\_T358K\_tmun0\_2p0 & numerical check & 0.607 & 1.510 & 117.3 & 285.3 & 11.01 & 5.14e-07 & 0.221 \\
\multicolumn{11}{@{}l}{\hspace{1em}\textit{elevated T; ATLAS tmu=1.5 applied (P-phonon as simulated)}} \\
\run{0055} & 2 & check\_overlap\_4um\_meshfixed\_dos384 & numerical check & 0.680 & 1.537 & 123.5 & 290.9 & 11.10 & 5.09e-07 & 0.043 \\
MEASURED\_6.3nm & 6.3 & measured & data & -- & -- & 124.5 & -- & -- & -- & -- \\
MEASURED\_13.2nm & 13.2 & measured & data & -- & -- & 136.1 & -- & -- & -- & -- \\
MEASURED\_31.8nm & 31.8 & measured & data & -- & -- & -- & -- & -- & -- & -- \\
MEASURED\_2nm & 2 & measured & data & -- & -- & 121.2 & -- & -- & -- & -- \\
\end{longtable}}


\section{Notes on non-standard runs}\label{sec:app-runs-notes}

\paragraph{\run{0010} (\SI{31.8}{nm}, back-sheet hypothesis A; status \texttt{NO\_EXPORT}).} The deck combined a
back-surface donor sheet of peak $3.23\times10^{14}\,\mathrm{cm^{-2}\,eV^{-1}}$ with $\Qf = 1.22\times10^{13}\cmsq$ to
describe the always-on \SI{31.8}{nm} device. ATLAS aborted in \texttt{SOLVE INIT} after \SI{17.9}{s}; no current was
written. The run is kept in the catalogue as a launch but contributes no number. The \SI{31.8}{nm} film was
subsequently excluded from the final set by the user.

\paragraph{\run{0011} (\SI{6.3}{nm}, stage 6, $\muband = 11.69$; status \texttt{NO\_EXECUTION\_JSON}).} The run was
part of a chain that was stopped before it executed; it has no \file{execution.json} and no currents. The value
$\muband = \SI{11.69}{cm^2/(V.s)}$ for \SI{6.3}{nm} was later used, together with $\Qf = 8.7\times10^{10}\cmsq$, in the
physical-structure run \run{0015}.

\paragraph{\run{0018} (\SI{13.2}{nm}, mesh $\times 0.7$; status \texttt{TIMEOUT}).} The refined mesh (24\,255 nodes)
exceeded the default timeout of \SI{900}{s}. The identical configuration was repeated with a timeout of \SI{1500}{s} as
\run{0028} (\SI{1068}{s}), which supplies check A$'$ of \cref{tab:num-convergence}.

\paragraph{\run{0027} (\SI{2}{nm}, contact overlap \SI{4}{\micro\metre}; MALFORMED, retracted).} The run is kept in the
catalogue, but its malformed mesh collapsed the drain, so its scored $\Ion$ change is not an overlap effect and the
overlap sensitivity is \emph{not determined} (full statement in \cref{sec:cal-param-control}).

\paragraph{Re-scored run \run{0012}.} The final \SI{2}{nm} run at DOS 96/48 was re-scored without a new launch after the
KCL acceptance rule was redefined on the median off-band current (check G$'$ of \cref{tab:num-convergence}); the original
\file{execution.json} is kept next to the re-scored one.

\paragraph{Rescaled curves \run{0039} and \run{0040}.} The \SI{6.3}{nm} and \SI{13.2}{nm} recalibration runs were
launched at $\muband = 11.41$ and \SI{60.39}{cm^2/(V.s)}. The exact recalibrated values, 11.58 and 61.60, follow from the
linearity of $\Id$ in MUN; wherever the report shows the final curves, the run's currents are multiplied by
$\times 1.015074$ and $\times 1.020113$, respectively, and the factor is stated. The catalogue lists the unscaled runs.

\paragraph{Temperature runs \run{0044}--\run{0047}.} These runs were planned with a temperature-independent band mobility,
but ATLAS applied its default constant-mobility exponent TMUN $= 1.5$ (\cref{sec:num-solver}). Their catalogue metrics are
the as-simulated values (variant P-phonon); the P-MTR variant is obtained by multiplying the currents by $(T/300)^{1.5}$
(\cref{ch:predictions}).

\section{History: the V0 and Astra packages}\label{sec:app-runs-history}

\paragraph{V0 (\file{IWO_ATLAS_Model_claude}, 34 launches).} The first package of this project fitted each thickness
separately with the same generator lineage, mesh, DOS discretization, interface regularization and solver controls as V1.
Its band mobilities were 16.8, 12.4, 57.6 and \SI{84}{cm^2/(V.s)} for 2.0, 6.3, 13.2 and \SI{31.8}{nm}; the
\SI{31.8}{nm} film required an additional back-surface donor sheet and a series resistance of about
$9\times10^{4}\,\Omega\,\mu$m. V0 contributes no calibrated curve to this report, but its numerical checks (mesh
$\times 0.5$ at \SI{2}{nm}, mesh $\times 0.7$ at \SI{13.2}{nm}, DOS 192/96, interface layer, tolerance, width
convention, reproducibility; documented in its \file{docs/VALIDATION_CHECKS.md}) apply to the V1 decks and are
included in \cref{tab:num-convergence}, always written as ``V0 run\_NNNN''. V0 also established two solver facts used
throughout: MOBILITY updates between SOLVE statements are ignored by ATLAS 5.28.1.R (V0 run\_0008), and the strict
continuity tolerance with XANDRNORM is needed in depletion (V0 run\_0006).

\paragraph{Astra (\file{astra_IWO_PRIORITY_THREE_20260912}, 70 solver invocations).} An independent earlier package
kept Conductor and Palladium volumes for the gate and the contacts and obtained the same class of fits. Its audit
documents (\file{audit/docs/SOURCES_AND_CORRECTIONS.md} and the \file{rebuild_20260912} notes) are used for three
cross-checks: electrons-only versus two-carrier solutions at \SI{2}{nm} (identical $\Id$; check I), repeated identical
runs (check H), and the thermal-generation bound on the off-current, 13 to 17 decades below the measured floors
(\cref{sec:inp-data}). The permittivities of \HfO{} (19.57) and IWO (9.30) from the SI of \citet{Januar2026} were also
recorded by the Astra audit. No Astra current is used as a calibrated or predicted curve in this report.

\begin{keyresult}
The V1 record contains 52 recorded launches, of which 49 produced scored currents (46 scored transfer curves and
3 $\Id$--$\Vd$ predictions), and 3 launches without
a usable curve (\run{0010} aborted, \run{0011} never executed, \run{0018} timed out and repeated as \run{0028}). One
scored run, \run{0027}, is malformed and retracted, so the contact-overlap sensitivity is not determined. Every simulated
number in the report refers to one of the catalogued runs; V0 (34 launches) and Astra (70 invocations) are used only as
documented cross-checks.
\end{keyresult}
